\begin{ex}
Trên cành cây có 2 chú chim sẻ đang đậu. Lát sau có thêm 1 chú chim chích và 2 chú chim khuyên cùng bay đến đậu chung cành. Hỏi lúc này trên cành cây có tất cả bao nhiêu chú chim?
\shortans{5}
\loigiai{
Ban đầu trên cành có 2 chú chim sẻ.\\
Số chim bay đến thêm là:
$$1 + 2 = 3\text{ (chú chim)}$$
Lúc này trên cành cây có tất cả số chú chim là:
$$2 + 3 = 5\text{ (chú chim)}$$
\textbf{Đáp số:} $5$.
}
\end{ex}

\begin{ex}
Trong một chiếc hộp có 5 viên bi gồm hai màu xanh và đỏ. Biết số viên bi xanh là 5 viên. Hỏi trong chiếc hộp đó có bao nhiêu viên bi đỏ?
\shortans{0}
\loigiai{
Trong hộp có tất cả 5 viên bi.\\
Vì số bi xanh đã là 5 viên nên số viên bi đỏ trong hộp là:
$$5 - 5 = 0\text{ (viên bi)}$$
\textbf{Đáp số:} $0$.
}
\end{ex}

\begin{ex}
Cho các số tự nhiên từ $0$ đến $5$: $0, 1, 2, 3, 4, 5$. Tìm số đứng liền sau của số lớn nhất trong các số bé hơn 4.
\shortans{4}
\loigiai{
Các số trong dãy bé hơn 4 là: $0, 1, 2, 3$.\\
Số lớn nhất trong các số này là số $3$.\\
Số đứng liền sau số $3$ trong dãy số tự nhiên là số $4$.\\
\textbf{Đáp số:} $4$.
}
\end{ex}

\begin{ex}
Tìm một số tự nhiên, biết rằng số đó đứng liền sau số 0 và đứng liền trước số 2 trong dãy số tự nhiên.
\shortans{1}
\loigiai{
Số đứng liền sau số 0 là số 1.\\
Số đứng liền trước số 2 là số 1.\\
Vậy số cần tìm thỏa mãn cả hai điều kiện trên là số $1$.\\
\textbf{Đáp số:} $1$.
}
\end{ex}

\begin{ex}
Bạn Nam có các tấm thẻ ghi các số từ 1 đến 5. Nam lấy ra hai tấm thẻ có tổng hai số bằng 5 và hiệu hai số bằng 1. Hỏi tấm thẻ có số lớn hơn ghi số mấy?
\shortans{3}
\loigiai{
Các cặp hai số có tổng bằng 5 là:
\begin{itemize}
    \item $1$ và $4$: Hiệu là $4 - 1 = 3$ (loại).
    \item $2$ và $3$: Hiệu là $3 - 2 = 1$ (chọn).
\end{itemize}
Vậy hai thẻ Nam lấy ra là thẻ số $2$ và thẻ số $3$.\\
Tấm thẻ có số lớn hơn ghi số $3$.\\
\textbf{Đáp số:} $3$.
}
\end{ex}

\begin{ex}
Có bao nhiêu số tự nhiên $x$ khác 0 thỏa mãn điều kiện: $x$ không vượt quá 5?
\shortans{5}
\loigiai{
Các số tự nhiên không vượt quá 5 gồm: $0, 1, 2, 3, 4, 5$.\\
Vì $x$ là số tự nhiên khác 0 nên $x$ có thể nhận các giá trị: $1, 2, 3, 4, 5$.\\
Đếm các số này, ta thấy có tất cả $5$ số.\\
\textbf{Đáp số:} $5$.
}
\end{ex}

\begin{ex}
Có ba bạn An, Bình và Cúc có số kẹo lần lượt là 1 cái, 4 cái và 5 cái. Hỏi bạn có nhiều kẹo nhất nhiều hơn bạn có ít kẹo nhất bao nhiêu cái kẹo?
\shortans{4}
\loigiai{
So sánh số kẹo của ba bạn: $1 < 4 < 5$.\\
Bạn có nhiều kẹo nhất là bạn Cúc ($5$ cái kẹo).\\
Bạn có ít kẹo nhất là bạn An ($1$ cái kẹo).\\
Số kẹo bạn Cúc nhiều hơn bạn An là:
$$5 - 1 = 4\text{ (cái kẹo)}$$
\textbf{Đáp số:} $4$.
}
\end{ex}

\begin{ex}
Trong giỏ có một số quả cam. Mẹ lấy ra 3 quả cho bà thì trong giỏ còn lại 2 quả cam. Hỏi ban đầu trong giỏ có bao nhiêu quả cam?
\shortans{5}
\loigiai{
Muốn tìm số quả cam ban đầu trong giỏ, ta lấy số cam mẹ đã lấy ra cộng với số cam còn lại trong giỏ:\\
Ban đầu trong giỏ có số quả cam là:
$$3 + 2 = 5\text{ (quả cam)}$$
\textbf{Đáp số:} $5$.
}
\end{ex}

\begin{ex}
Cho dãy số tăng dần: $1, 3, 5$. Hỏi số đứng ở vị trí thứ hai của dãy số là số nào?
\shortans{3}
\loigiai{
Dãy số đã cho gồm 3 số:
\begin{itemize}
    \item Số đứng ở vị trí thứ nhất là $1$.
    \item Số đứng ở vị trí thứ hai là $3$.
    \item Số đứng ở vị trí thứ ba là $5$.
\end{itemize}
Vậy số đứng ở vị trí thứ hai của dãy số là số $3$.\\
\textbf{Đáp số:} $3$.
}
\end{ex}

\begin{ex}
Một hàng dọc gồm có 5 bạn học sinh. Bạn Hoa đứng ở vị trí thứ 3 tính từ đầu hàng xuống. Hỏi phía sau bạn Hoa có bao nhiêu bạn học sinh?
\shortans{2}
\loigiai{
Hàng dọc có tất cả 5 bạn.\\
Tính từ đầu hàng đến bạn Hoa có 3 bạn (kể cả Hoa).\\
Số bạn học sinh đứng ở phía sau bạn Hoa là:
$$5 - 3 = 2\text{ (bạn)}$$
\textbf{Đáp số:} $2$.
}
\end{ex}

\begin{ex}
Cho dãy số tự nhiên liên tiếp: $0, 1, 2, \dots, 4, 5$. Lấy số lớn nhất trong dãy trừ đi số còn thiếu ở vị trí dấu ba chấm thì được kết quả bằng bao nhiêu?
\shortans{2}
\loigiai{
Dãy số tự nhiên liên tiếp từ 0 đến 5 là: $0, 1, 2, 3, 4, 5$.\\
Số còn thiếu ở vị trí dấu ba chấm là số $3$.\\
Số lớn nhất trong dãy là số $5$.\\
Lấy số lớn nhất trừ đi số còn thiếu, ta được:
$$5 - 3 = 2$$
\textbf{Đáp số:} $2$.
}
\end{ex}

\begin{ex}
Bạn Minh đếm lùi từ 5 về 0: $5, 4, \dots, 2, 1, 0$. Hỏi số đứng ngay trước số 2 trong lượt đếm của Minh là số nào?
\shortans{3}
\loigiai{
Dãy số đếm lùi đầy đủ từ 5 về 0 là: $5, 4, 3, 2, 1, 0$.\\
Trong lượt đếm lùi này, số được gọi ngay trước số 2 là số $3$.\\
\textbf{Đáp số:} $3$.
}
\end{ex}

\begin{ex}
Lan có 2 cái kẹo, mẹ cho Lan thêm 1 cái kẹo. Sau đó Lan cho em 1 cái kẹo. Hỏi lúc này Lan còn lại bao nhiêu cái kẹo?
\shortans{2}
\loigiai{
Sau khi mẹ cho thêm, Lan có tất cả số kẹo là:
$$2 + 1 = 3\text{ (cái kẹo)}$$
Sau khi Lan cho em 1 cái kẹo, Lan còn lại số kẹo là:
$$3 - 1 = 2\text{ (cái kẹo)}$$
\textbf{Đáp số:} $2$.
}
\end{ex}

\begin{ex}
Trên bàn có một đĩa đựng 4 quả táo và một đĩa đựng 1 quả lê. Nam ăn hết 2 quả táo. Hỏi trên bàn lúc này còn lại tất cả bao nhiêu quả của cả hai loại?
\shortans{3}
\loigiai{
Số quả táo còn lại sau khi Nam ăn là:
$$4 - 2 = 2\text{ (quả táo)}$$
Số quả của cả hai loại còn lại trên bàn là:
$$2 + 1 = 3\text{ (quả)}$$
\textbf{Đáp số:} $3$.
}
\end{ex}

\begin{ex}
Trong các số $0, 1, 2, 3, 4, 5$, có bao nhiêu số vừa lớn hơn 1 vừa bé hơn 5?
\shortans{3}
\loigiai{
Các số thỏa mãn đồng thời lớn hơn 1 và bé hơn 5 là: $2, 3, 4$.\\
Đếm các số tìm được, ta có tất cả $3$ số.\\
\textbf{Đáp số:} $3$.
}
\end{ex}

\begin{ex}
Tìm số tự nhiên $a$ biết rằng khi lấy $a$ cộng thêm 2 thì được kết quả bằng 5.
\shortans{3}
\loigiai{
Theo đề bài, ta có phép tính:
$$a + 2 = 5$$
Muốn tìm số hạng chưa biết $a$, ta lấy tổng trừ đi số hạng đã biết:
$$a = 5 - 2 = 3$$
Vậy $a = 3$.\\
\textbf{Đáp số:} $3$.
}
\end{ex}

\begin{ex}
Một chú mèo có 4 cái chân. Khi chú mèo co 1 chân lên thì có bao nhiêu chân của chú mèo đang chạm đất?
\shortans{3}
\loigiai{
Một chú mèo có 4 cái chân.\\
Khi chú mèo co 1 chân lên, số chân còn lại chạm đất là:
$$4 - 1 = 3\text{ (chân)}$$
\textbf{Đáp số:} $3$.
}
\end{ex}

\begin{ex}
Bạn Hùng có một chiếc xe đạp 2 bánh và một chiếc xe ô tô con 4 bánh. Hỏi chiếc ô tô con có nhiều hơn chiếc xe đạp bao nhiêu bánh xe?
\shortans{2}
\loigiai{
Chiếc xe đạp có 2 bánh xe.\\
Chiếc xe ô tô con có 4 bánh xe.\\
Số bánh xe của chiếc ô tô con nhiều hơn chiếc xe đạp là:
$$4 - 2 = 2\text{ (bánh xe)}$$
\textbf{Đáp số:} $2$.
}
\end{ex}

\begin{ex}
Có 5 que tính được xếp thành một hàng ngang. Em nhấc bớt đi 2 que tính ở hai đầu hàng. Hỏi trong hàng còn lại bao nhiêu que tính?
\shortans{3}
\loigiai{
Ban đầu có tất cả 5 que tính.\\
Sau khi nhấc đi 2 que tính ở hai đầu, số que tính còn lại trong hàng là:
$$5 - 2 = 3\text{ (que tính)}$$
\textbf{Đáp số:} $3$.
}
\end{ex}

\begin{ex}
Cho các số: $0, 2, 4, 1, 5$. Trong các số đã cho, có bao nhiêu số lớn hơn số 2?
\shortans{2}
\loigiai{
So sánh từng số đã cho với số 2:
\begin{itemize}
    \item Các số bé hơn 2 là: $0, 1$.
    \item Số bằng 2 là: $2$.
    \item Các số lớn hơn 2 là: $4, 5$.
\end{itemize}
Vậy có tất cả $2$ số lớn hơn số 2.\\
\textbf{Đáp số:} $2$.
}
\end{ex}

\begin{ex}
Cho các số: $5, 0, 3, 2, 4$. Sắp xếp các số này theo thứ tự từ bé đến lớn, hỏi số đứng ở vị trí chính giữa (thứ ba) là số nào?
\shortans{3}
\loigiai{
Sắp xếp các số theo thứ tự từ bé đến lớn, ta được:
$$0 < 2 < 3 < 4 < 5$$
Dãy số gồm 5 số, số đứng ở vị trí chính giữa (vị trí thứ ba) là số $3$.\\
\textbf{Đáp số:} $3$.
}
\end{ex}

\begin{ex}
Lấy số lớn nhất trong các số từ 0 đến 5 trừ đi số bé nhất trong các số này, ta được kết quả là bao nhiêu?
\shortans{5}
\loigiai{
Trong các số từ 0 đến 5:
\begin{itemize}
    \item Số lớn nhất là số $5$.
    \item Số bé nhất là số $0$.
\end{itemize}
Hiệu giữa số lớn nhất và số bé nhất là:
$$5 - 0 = 5$$
\textbf{Đáp số:} $5$.
}
\end{ex}

\begin{ex}
Tìm số đứng liền trước của số lớn nhất có một chữ số trong phạm vi từ 0 đến 5.
\shortans{4}
\loigiai{
Trong phạm vi từ 0 đến 5, số lớn nhất là số $5$.\\
Số đứng liền trước số 5 trong dãy số tự nhiên là số $4$ (vì $5 - 1 = 4$).\\
\textbf{Đáp số:} $4$.
}
\end{ex}

\begin{ex}
Mai có 5 bông hoa. Mai cắm vào lọ màu xanh nhiều hơn lọ màu đỏ 1 bông hoa và cắm hết số hoa. Hỏi lọ màu xanh có bao nhiêu bông hoa?
\shortans{3}
\loigiai{
Tổng số hoa cắm vào hai lọ là 5 bông.\\
Ta tìm hai số có tổng bằng 5 và số hoa ở lọ xanh hơn lọ đỏ 1 bông:
$$2 + 3 = 5\text{ và } 3 - 2 = 1$$
Vậy lọ màu đỏ có 2 bông hoa và lọ màu xanh có 3 bông hoa.\\
\textbf{Đáp số:} $3$.
}
\end{ex}

\begin{ex}
Bình có 3 viên bi xanh và 2 viên bi đỏ. Bạn An có số viên bi đúng bằng số bi đỏ của Bình. Hỏi bạn Bình có nhiều hơn bạn An bao nhiêu viên bi?
\shortans{3}
\loigiai{
Tổng số viên bi của Bình là:
$$3 + 2 = 5\text{ (viên bi)}$$
Số viên bi của An bằng số bi đỏ của Bình nên An có 2 viên bi.\\
Số viên bi Bình nhiều hơn An là:
$$5 - 2 = 3\text{ (viên bi)}$$
\textbf{Đáp số:} $3$.
}
\end{ex}
